% Author: Animesh Garg
% Date: Aug 31 2022

% ---------------------------------
% Common base packages
% ---------------------------------

% \usepackage{etoolbox} % fix space between abstract header and abstract body
% \usepackage[utf8]{inputenc} % allow utf-8 input
% \usepackage[T1]{fontenc}    % use 8-bit T1 fonts
% \usepackage[english]{babel} %for hyphenation rules
% % \usepackage{microtype}      % microtypography
% \usepackage[protrusion=true,expansion=true]{microtype}
% %\usepackage{lmodern}
\usepackage{comment}
% \usepackage{cancel}
% \usepackage{csquotes}
% \usepackage{listings}
% \usepackage{nameref}
% % \usepackage{slashbox}
% \usepackage{paralist}
% \usepackage{xspace}
% \usepackage{setspace}
% % \usepackage{listings} for text to reproduce verbatim
% \usepackage{makeidx}
% \usepackage{pdfpages}

% ---------------------------------
% MATH & Equations
% ---------------------------------
\usepackage{amsmath,amssymb} % define this before the line numbering.
% \usepackage{amsfonts}       % blackboard math symbols
% \usepackage{mathptmx}

% \usepackage{newtxtext,newtxmath}

\usepackage{mathtools}
\usepackage{array} %for table entries to be in center of cell
\usepackage{nicefrac}       % compact symbols for 1/2, etc.
\usepackage{breqn}          % dmath multiline
\usepackage{textcomp}
\usepackage{bm} %bold symbols in math mode
\usepackage{pifont}
\usepackage[
  separate-uncertainty = true,
  multi-part-units = repeat
]{siunitx}

% ---------------------------------
% Figures & Floats
% ---------------------------------
\usepackage{graphicx}
% \usepackage[pdftex]{graphicx}
\usepackage{adjustbox} %add additional box operations in addition to resizebox
\usepackage{epsfig}
% \pdfoutput=1
% \graphicspath{{./figures/}}
% \DeclareGraphicsExtensions{.pdf,.jpeg,.png,.jpg}

\usepackage{float}
\usepackage{subcaption}
\usepackage[font=small,labelfont=bf]{caption}
%\usepackage[justification=raggedright]{caption}	% makes captions ragged right 
\usepackage{lscape}      %Useful for wide tables or figures.
%\usepackage{subfig}
%\usepackage{subfigure}		% For subfigures

\usepackage{tikz}  %make graphs
\usepackage{makecell}
\usepackage{overpic}
% pdfTeX-only setting; Tectonic/XeTeX selects the PDF version via its driver.
\ifdefined\pdfminorversion
  \pdfminorversion=4
\fi

% ---------------------------------
% Algorithms
% ---------------------------------
\usepackage{algorithm}
\usepackage[noend]{algpseudocode}

% ---------------------------------
%Table Stuff
% ---------------------------------
\usepackage{array} %for table entries to be in center of cell
\usepackage{tabularx}
\usepackage{multirow}
\usepackage{multicol}
\usepackage{booktabs}   % Publication quality tables
% \usepackage{colortbl} %dont load with xcolor, loaded in colors
\usepackage{tablefootnote}

% ---------------------------------
% LISTS compact itemize/enumeration (e.g. contributions)
% ---------------------------------
\usepackage{paralist}
\let\labelindent\relax % Avoid collision with the IEEE class definition.
\usepackage{enumitem}
\setlist[itemize]{noitemsep,leftmargin=*,topsep=0in}
\setlist[enumerate]{noitemsep,leftmargin=*,topsep=0in}

% ---------------------------------
% URLS, Colors, Links
% ---------------------------------
\usepackage{url}            % simple URL typesetting
% \usepackage[table]{xcolor}
\PassOptionsToPackage{table}{xcolor}
% \usepackage{color}
\usepackage{color,colortbl} %superceded by xcolor
\usepackage{soul}

% \usepackage[capitalize, noabbrev]{cleveref}
\PassOptionsToPackage{capitalize, noabbrev}{cleveref}
\usepackage[pagebackref=false,breaklinks=true,colorlinks,urlcolor=blue,citecolor=blue,linkcolor=blue,bookmarks=false]{hyperref}

% ---------------------------------
% Citation
% ---------------------------------
% \makeatletter
% \let\NAT@parse\undefined
% \makeatother
% \usepackage[numbers,sort&compress]{natbib}

% ---------------------------------
% USER defined macros: shortcuts for standard references
% ---------------------------------

\newcommand{\Fig}[1]{Fig.~\ref{fig:#1}}
\newcommand{\Figure}[1]{Figure~\ref{fig:#1}}
\newcommand{\Tab}[1]{Tab.~\ref{tab:#1}}
\newcommand{\Table}[1]{Table~\ref{tab:#1}}
\newcommand{\eq}[1]{(\ref{eq:#1})}
\newcommand{\Eq}[1]{Eq.~\ref{eq:#1}}
\newcommand{\Equation}[1]{Equation~\ref{eq:#1}}
\newcommand{\Sec}[1]{Sec.~\ref{sec:#1}}
\newcommand{\Section}[1]{Section~\ref{sec:#1}}
\newcommand{\Appendix}[1]{Appendix~\ref{app:#1}}

% ---------------------------------
% USER defined macros: lorem (i.e. filler latin text)
% ---------------------------------
\usepackage{blindtext}
\usepackage{bbm}
\newcommand{\lorem}[1]{\todo{\blindtext[#1]}}
\usepackage{lipsum}
\newcommand{\loremnew}[1]{\todo{\lipsum[1][0-#1]}}

% ---------------------------------
% USER defined macros: footnote without a number
% ---------------------------------
\newcommand\blfootnote[1]{%
  \begingroup
  \renewcommand\thefootnote{}\footnote{#1}%
  \addtocounter{footnote}{-1}%
  \endgroup
}

% ---------------------------------
% USER defined macros: Space Tweaks
% ---------------------------------
\setlength{\abovecaptionskip}{1.5mm}
\setlength{\belowcaptionskip}{1.0mm} 
\setlength{\textfloatsep}{1.5mm}
\setlength{\dbltextfloatsep}{1.5mm}

%\renewcommand{\baselinestretch}{.9945} %space between lines - use as last resort
\renewcommand{\arraystretch}{1.1}

% Spacing before and after section headings
\usepackage{titlesec}
\titlespacing{\section}{0pt}{0.3\baselineskip}{0.25\baselineskip}
\titlespacing{\subsection}{0pt}{0.2\baselineskip}{0.15\baselineskip}
\titlespacing{\subsubsection}{0pt}{0.05\baselineskip}{0.03\baselineskip}

%% paragraph (fine tune spacing close to deadline)
\renewcommand{\paragraph}[1]{\vspace{0.2em}\noindent\textit{#1} --}
% \newcommand\numberthis{\addtocounter{equation}{1}\tag{\theequation}}

% ---------------------------------
% USER defined macros: Comments/edits
% ---------------------------------
\newcommand{\todo}[2]{{\color{orange}#1:#2}} %< I changed something and I want you to see it
\newcommand{\TODO}[2]{\sethlcolor{yellow}\hl{[#1: #2]}}%< inlined comment for max visibility
\newcommand{\todofoot}[2]{\footnote{{\color{red}#1:#2}}}%< useful for ~deadline (no layout changes)

%%
%% Colors
%%
% \definecolor{Lightapricot}{rgb}{0.99,0.84,0.69}



%%%%%%%%%%%%%%%%%%%%%%%%%%%%%%%%
% New list style
%%%%%%%%%%%%%%%%%%%%%%%%%%%%%%%%

\usepackage{listings}
\lstset{escapeinside={*@}{@*},
breaklines=true, breakatwhitespace=false,}
\def\code#1{\texttt{#1}}
% https://www.overleaf.com/learn/latex/Code_listing
\definecolor{codegreen}{rgb}{0,0.6,0}
\definecolor{codegray}{rgb}{0.4,0.4,0.4}
\definecolor{codepurple}{rgb}{0.5,0,0.9}
\definecolor{backcolour}{rgb}{0.95,0.95,0.95}

\definecolor{lightgray}{rgb}{0.9,0.9,0.9}
\definecolor{lightpink}{rgb}{0.98,0.85,0.86}
\definecolor{lightblue}{rgb}{0.68,0.84,0.9}

\lstdefinestyle{mystyle}{
    backgroundcolor=\color{backcolour},   
    commentstyle=\color{codegreen},
    keywordstyle=\color{magenta},
    numberstyle=\tiny\color{codegray},
    stringstyle=\color{codepurple},
    basicstyle=\fontsize{8}{8}\selectfont\ttfamily\ttfamily,
    breakatwhitespace=false,         
    breaklines=true,
    breakindent=0pt,
    captionpos=b,                    
    keepspaces=true,                 
    numbers=none,                    
    numbersep=5pt,                  
    showspaces=false,                
    showstringspaces=false,
    showtabs=false,                  
    tabsize=2
}
\lstset{style=mystyle}
\renewcommand{\lstlistingname}{Snippet}
